El resumen de una tesis es una síntesis breve y precisa que presenta los aspectos fundamentales de la investigación. Su finalidad es ofrecer al lector una visión general del trabajo, permitiéndole identificar rápidamente el objetivo, la metodología empleada, los resultados más relevantes y las conclusiones alcanzadas. Un buen resumen debe reflejar fielmente el contenido de la tesis, destacar sus principales aportaciones y comunicar la información de manera clara, objetiva y comprensible. Asimismo, debe redactarse en un único párrafo, sin incluir citas bibliográficas, figuras, tablas ni información ajena al estudio. Para garantizar su concisión y utilidad, su extensión no deberá exceder las 200 palabras, permitiendo así una comprensión rápida y completa de la investigación realizada.
